\chapter{결론 및 향후 연구 방향}
\section{결론}
SO-101의 CV+IK 기반 블록 이동 시스템을 구현하고 실기에서 발생한 인식 혼동, 좌표 편차와 빈손 파지에 대응하였다. 색상·형상 기반 대상 선택, 중립 손목 방향과 대체 파지 후보, 그리퍼 위치·부하 확인, 재시도와 슬롯 배치를 하나의 상태 흐름으로 연결하였다. 카메라 웹 뷰어와 캘리브레이션 도구에는 영상, 대응점과 관절 상태를 함께 기록하여 실패 단계를 추적하였다.

초기 180초 제한 시험에서는 파지 후보를 14회 실행하였다. 센서 판정은 HELD 4회, EMPTY 9회였으며 1회는 판정 전에 중단되었다. 노란 블록은 같은 PICK의 세 번째 후보에서 파지했고, 빨간 블록은 재선택 뒤 파지하여 운반과 슬롯 해제를 마쳤다. 단일 블록 시험에서는 파지와 운반에 성공했지만 위치를 바꾼 뒤에는 빈손 파지가 발생했다.

초기 시험은 네 번의 슬롯 해제 후 마지막 블록이 남아 미완료로 판정하였다. 빨간 블록의 첫 PICK에는 51.5초가 걸려 한 블록에 여러 후보를 실행할 때의 시간 부담이 드러났다. 당시 좌표 정합 오차도 RMS 8.18\,mm, 최대 LOO 26.64\,mm로 개발 목표를 넘었다. ACT와의 성능 차이는 같은 조건에서 측정하지 못했다.

로봇팔 여섯 축의 나사를 다시 조여 백래시를 줄이고, Z축 높이와 손목 회전각을 일정하게 유지하며 캘리브레이션을 다시 수행하였다. 최종 RMS는 5.1454\,mm, 최대 LOO 오차는 13.5348\,mm로 줄었다. RMS는 5\,mm 기준에 근소하게 미달했지만 중간 결과보다 37.1\% 개선되었다. ArUco 마커와 손목 카메라 장착에 이어 그리퍼 회전 재시도와 하강 민감도 조정도 적용하였다. 이후 1차 미션은 28회 모두 180초 안에 완료하였다. 같은 조건의 2차 미션에서 3단 적재는 시험한 범위에서 모두 성공했다. 4단 적재는 추가 조정 후 1회 시도에 성공하여 5초 유지까지 확인하였다. 앞으로는 1차 미션의 위치·방향별 지표를 수집하고, 4단 적재를 반복 평가한 뒤 5단 적재를 시험할 예정이다.

\section{향후 연구 방향}
\subsection{동일 조건 반복 평가와 시간 예산 개선}
그리퍼 회전 재시도와 하강 민감도 조정 후 180초 제한 시험을 28회 반복해 모두 완료하였다. 다만 인정 블록 수, 후보별 센서 판정과 소요 시간을 세분화해 기록하지 않아 완료 여부 외의 지표는 남아 있지 않다.

후속 시험에서는 후보별 시간 상한과 다음 블록으로 넘어갈 시점을 정하고, 위치·방향별 파지 성공률을 측정해야 한다. 무작위 배치에 대한 반복 평가도 필요하다. 실제 턱과 블록의 상대 위치 측정, 파지 직전 영상 보정, 검출 누락과 물리적 완료 조건의 차이도 함께 검토한다.

\subsection{적재 높이별 동작 개선과 실기 평가}
\label{sec:stack-extension}
정확도 개선 후 타워 접근과 하강 자세를 등록하여 3단과 4단 적재를 시험하였다. 3단 적재는 시험한 범위에서 모두 성공했다. 4단 적재는 처음 8회 중 1회 성공했으며, 그리퍼 회전 재시도와 하강 민감도 조정 후에는 추가 1회도 성공했다. 총 9회 중 2회로 표본이 적어 반복 성공률로 해석하기 어렵다. 접촉 감지와 되돌림 동작이 타워를 흔드는 정도, 실패 단계별 원인도 분석해야 한다. 4단의 반복성을 확인한 뒤 아직 시도하지 않은 5단 적재로 확장할 계획이다. 유격 정비, ArUco 마커, 손목 카메라, 재시도 로직과 하강 민감도의 기여도는 분리해서 측정하지 못했다.

평면 변환은 블록 한 개의 윗면을 기준으로 하므로 타워가 높아졌을 때 같은 변환을 그대로 적용할 수 없다. 초기에는 고정한 타워 위치와 기록 자세를 사용하고, 높이 변화에 따른 관측 오차와 적재물 재선택 문제를 별도로 평가한다.

\subsection{규칙 기반 실행의 학습 데이터 활용}
\label{sec:dataset-extension}
Task3은 CV+IK 실행 중 영상, 관절 상태와 전달 명령을 같은 시간 기준으로 기록한다(\ref{sec:task3-results}절). 사람이 블록을 재배치하면 로봇이 home부터 배치까지 한 사이클을 반복하고, 성공한 사이클만 LeRobotDataset 형식으로 저장한다. 첫 실행에서 5에피소드를 수집해 파일 구조와 프레임 정합성을 확인하였다.

데이터에는 더 많은 배치와 접근 방향이 고르게 포함되어야 한다. 현재 에피소드는 파지부터 배치까지의 전체 사이클이므로 ACT 학습 범위를 파지로 한정하려면 구간 분리(segmentation) 기준도 필요하다. 최초 파지 실패 시 에피소드를 폐기하기 때문에 실패·복구 시연이 없고, 2차 미션의 접촉 기반 적재 정렬도 포함하지 않았다. 제어 주기가 목표보다 26\% 느린 원인을 분석하고 ArUco 기반 실측 좌표를 함께 기록하는 방안도 검토할 예정이다.

LeRobotDataset의 영상, 상태, 행동과 에피소드 구성을 참고하되 사용할 버전과의 호환성을 확인한다\cite{lerobot-dataset}. 모방학습이나 SmolVLA에 활용하려면 측정 관절값과 전달 명령을 구분하고 작업 지시 및 센서 라벨의 신뢰도를 검토해야 한다\cite{smolvla}. 강화학습용으로는 다음 관측, 보상과 종료 조건을 추가 정의한다. 파지 판정과 배치 성공은 별도 라벨로 관리하며, 유효 에피소드 수와 기록 누락, 사람의 작업시간으로 수집 방식을 평가할 계획이다. 이 데이터로 ACT를 재학습하거나 성능을 비교하는 단계는 아직 수행하지 못하였다.
\FloatBarrier
